\chapter{آماده‌سازی پایتون و پروژه}

\begin{goalbox}
یک محیط جداگانهٔ پایتون بسازید و تنها بستهٔ لازم را نصب کنید.
\end{goalbox}

پایتون نسخهٔ ۳٫۱۰ یا جدیدتر را نصب کنید. نسخه را در ترمینال، پاورشل یا ترمینال
داخل ویرایشگر بررسی کنید:

\begin{LTR}
\begin{lstlisting}[language=bash,numbers=none]
python --version
\end{lstlisting}
\end{LTR}

در بعضی سیستم‌های مک یا لینوکس باید از دستور
\lr{python3}
استفاده کنید. نتیجه‌ای مانند
\lr{Python 3.12.4}
مناسب است.

\section{بازکردن پوشهٔ کد}

مخزن را دریافت کنید و وارد پوشهٔ کد شوید:

\begin{LTR}
\begin{lstlisting}[language=bash,numbers=none]
git clone https://github.com/alirezaaies/telbot-simple-bot.git
cd telbot-simple-bot/code
\end{lstlisting}
\end{LTR}

اگر فایل فشرده را دانلود کرده‌اید، آن را استخراج کنید و ترمینال را در پوشهٔ
\lr{code}
باز کنید.

\section{ساخت محیط جداگانه}

محیط مجازی، بستهٔ این آموزش را از پروژه‌های دیگر پایتون جدا نگه می‌دارد.

در مک یا لینوکس:
\begin{LTR}
\begin{lstlisting}[language=bash,numbers=none]
python3 -m venv .venv
source .venv/bin/activate
\end{lstlisting}
\end{LTR}

در پاورشل ویندوز:
\begin{LTR}
\begin{lstlisting}[language={},numbers=none]
python -m venv .venv
.venv\Scripts\Activate.ps1
\end{lstlisting}
\end{LTR}

ترمینال معمولاً عبارت
\lr{(.venv)}
را پیش از نشانگر خود نمایش می‌دهد. سپس وابستگی را نصب کنید:

\begin{LTR}
\begin{lstlisting}[language=bash,numbers=none]
python -m pip install --upgrade pip
python -m pip install -r requirements.txt
\end{lstlisting}
\end{LTR}

فایل نیازمندی‌ها فقط شامل بستهٔ
\lr{pyTelegramBotAPI}
است. نام آن هنگام واردکردن در پایتون
\lr{telebot}
است؛ این تفاوت طبیعی است.

\begin{checkpointbox}
دستور زیر را اجرا کنید. اگر عبارت
\lr{telebot is ready}
نمایش داده شد، نصب کامل است.
\begin{LTR}
\begin{lstlisting}[language=bash,numbers=none]
python -c "import telebot; print('telebot is ready')"
\end{lstlisting}
\end{LTR}
\end{checkpointbox}
